\section{An obstructed repeated-factor application}\label{sec:application}

\subsection{The geometric input}
We use the stable bundle constructed in
\cite[Section~6]{wuebben15}. We recall its geometry
and deformation equations to specify the entire
background of the application.
On a base $\PP^2$, choose smooth cubic curves
defined by $R,P,Q$, meeting pairwise transversely
and with no common point. Fix $p_0\in\{P=Q=0\}$.
In $\PP^2_{[x:y:z]}\times\PP^2$ set
\begin{equation}\label{eq:bicubic}
 X=\{x^3R+y^3P+z^3Q=0\},\qquad
 C=\{x=0\}\times\{p_0\}.
\end{equation}
The threefold is smooth: if all fiber coordinates
are nonzero, the fiber derivatives would force
$R=P=Q=0$; with only one or two nonzero coordinates,
smoothness or transversality of the corresponding
base cubics prevents the total differential from
vanishing. Adjunction and the restriction sequence
give $K_X\simeq\OO_X$ and
$H^1(X,\OO_X)=H^2(X,\OO_X)=0$.

The Serre construction \cite{arrondo}, used for obstructed
bundles by Thomas \cite{thomas}, gives a locally free rank-two
bundle
\begin{equation}\label{eq:serre-bundle}
 0\longrightarrow\OO_X\longrightarrow E
       \longrightarrow I_C\OO_X(-2,2)\longrightarrow0,
 \qquad A=E(1,-1),\qquad \det A=\OO_X.
\end{equation}
Let $\omega_1,\omega_2$ denote the restricted hyperplane
classes. The geometric input supplies stability of $A$
for $N\omega_1+\omega_2$, $1<N<1+\sqrt3$, and the full
fixed-determinant analytic deformation ring
\begin{equation}\label{eq:scalar-germ}
 \C\{a,b\}/(a^2b,ab^2).
\end{equation}
In particular, $\dim\Ext^1(A,A)=2$, the quadratic
obstruction vanishes, and the cubic obstruction is
essential.

The calculations behind these assertions are concrete.
With $U=Q/R$ and $V=P/R$ near $p_0$, the local equation
is $x^3+Vy^3+Uz^3=0$. The full ambient line chart
$x=ay+bz$, $U=c$, $V=d$ gives coefficients
\[
 a^3+d,\qquad 3a^2b,\qquad 3ab^2,\qquad b^3+c.
\]
Eliminating $c,d$ gives \eqref{eq:scalar-germ}.
The normal bundle is $\OO_C(1)\oplus\OO_C(-3)$.
The local Serre generator is nowhere zero, and
$\OO_X(2,-2)$ is acyclic, proving local freeness
and the required global extension. The vanishings
$h^0(E)=1$ and $H^1(E)=0$ identify the curve and
bundle deformation functors, including their
nonreduced analytic structure: the section lifts
uniquely up to a unit, and its regular zero scheme
is the inverse construction. For stability, the
twists of $E$ with sections lie in the three upper
quadrants based at $(0,0),(2,-1),(3,-2)$.
Their first-Chern degrees against
$(N\omega_1+\omega_2)^2$ are respectively
$6(N^2-1)$, $6(2N+1)$ and $6(2+2N-N^2)$,
all positive on the stated interval.
These are the established geometric inputs;
the dynamical application below uses the full
deformation problem.

Fix a K\"ahler form in $2\omega_1+\omega_2$ and
the normalized HE metric of $A$, and put
$E_0=A\otimes\C^3$. Then
\begin{equation}\label{eq:application-data}
 G=\mathrm{SL}(3,\C),\quad K=\mathrm{SU}(3),\quad
 \dim_{\R}\ha=8,\qquad
 V=\Ext^1(A,A)\otimes\End\C^3,\quad \dim_{\C}V=18.
\end{equation}
The determinant condition removes no degree-one
harmonic direction because $H^1(X,\OO_X)=0$.
There is one distinct stable factor, so $b_p=0$
for every nearby polarization.

\begin{corollary}\label{cor:bicubic}
Theorem~\ref{thm:main} applies to every integrable
point of the full eighteen-coordinate deformation
problem at $A\otimes\C^3$, with common coefficient
and normal radii and its stated confinement hypothesis.
Its phase and relative tracking constants remain
uniform through all stabilizer changes and arbitrarily
separated positive finite scales.
\end{corollary}

\begin{proof}
The construction verifies a fixed smooth compact
K\"ahler background, local freeness, stability of
the factor and the determinant normalization.
All multiplicity directions in \eqref{eq:application-data}
are retained. Theorem~\ref{thm:main} imposes no
restriction on the order of the obstruction or
the rank of $C_x$, so applies to the entire zero
locus of the actual canonical $\kappa_p$.
\end{proof}

The matrix obstruction in this statement is defined
by \eqref{eq:kuranishi}. It is not obtained by
substituting matrices into the scalar equations
\eqref{eq:scalar-germ}. On the scalar subspace
$x=u\otimes I_3$, the canonical recursion retains
the cubic obstruction. A quadratic truncation
would therefore admit false integrable points.

\subsection{Noncommuting integrable coefficients}
The Calabi--Yau duality supplies a useful additional
fact for this two-dimensional deformation space.

\begin{lemma}\label{lem:cup}
Every product of two elements of $\Ext^1(A,A)$
vanishes in $\Ext^2(A,A)$.
\end{lemma}

\begin{proof}
Take a basis $a,b$ of $\Ext^1(A,A)$. The zero
quadratic obstruction gives $a^2=b^2=0$ and
$ab+ba=0$. Calabi--Yau trivialization and the
Serre trace give a perfect pairing
\[
 \Ext^2(A,A)\times\Ext^1(A,A)\longrightarrow\C,
             \qquad (u,v)\longmapsto\Tr(uv).
\]
Graded cyclicity gives
$\Tr(aba)=\Tr(aab)=0$: moving the last degree-one
factor past the first two has sign
$(-1)^{1\cdot2}=+1$. Also $\Tr(abb)=0$.
Thus $ab$ pairs to zero with the basis $a,b$,
and perfection implies $ab=0$. Then $ba=0$,
and bilinearity proves the result.
\end{proof}

For harmonic representatives $\alpha,\beta$,
the lemma implies
$H_2(\alpha\wedge\beta)=H_2(\beta\wedge\alpha)=0$,
as well as the square identities. The trace in
the proof is the Serre trace, or integration of
the bundle trace against the holomorphic volume
form. Zero scalar quadratic obstruction alone
would not imply this product vanishing in an
arbitrary deformation problem.

Choose $\alpha,\beta$ orthonormal in harmonic
$L^2$. For small complex parameters $s,t$, set
\begin{equation}\label{eq:noncommuting-x}
 x_{s,t}=s\alpha\otimes E_{12}+t\beta\otimes E_{23},
 \qquad B=d^*G_2(\alpha\wedge\beta).
\end{equation}

\begin{proposition}\label{prop:noncommuting}
The exact canonical reconstruction is
\begin{equation}\label{eq:noncommuting-a}
 a(x_{s,t})=s\alpha\otimes E_{12}+t\beta\otimes E_{23}
                         -stB\otimes E_{13},\qquad
 \kappa(x_{s,t})=0.
\end{equation}
For $st\ne0$ the harmonic matrices do not commute.
The resulting holomorphic bundle is locally free
with fixed determinant and a three-step filtration
whose successive quotients are $A$.
\end{proposition}

\begin{proof}
For the displayed $a$, the only nonzero product
in $a\wedge a$ is
$st(\alpha\wedge\beta)\otimes E_{13}$.
Indeed $E_{12}E_{23}=E_{13}$, whereas the reverse
product, both squares, and every product involving
the correction $E_{13}$ vanish. Thus the displayed
form solves the full equation
$a=x-d^*G_2(a\wedge a)$ exactly. Contraction
uniqueness identifies it with the canonical
reconstruction, and Lemma~\ref{lem:cup} gives its
harmonic obstruction zero. The full Kuranishi
construction gives actual integrability.

Strict upper triangularity preserves the standard
flag, with diagonal quotients $A$. The integrable
smooth Dolbeault realization is locally free, and
trace zero fixes the determinant. Finally
$[sE_{12},tE_{23}]=stE_{13}\ne0$. A nonzero
commutator cannot become zero under simultaneous
conjugation, so this family is outside every
conjugate of a commuting or diagonal tuple.
\end{proof}

These bundles are semistable, being successive
extensions of equal-slope stable bundles. This
does not assert existence of an HE metric on a
nonsplit member. The example solves the complete
recursion on a particular noncommuting family;
it does not give all matrix obstruction equations.

\subsection{Changing stabilizers and separated scales}
Write $u=|s|^2$, $v=|t|^2$. On
$\mathfrak{sl}_3(\C)$ with Hilbert--Schmidt mass,
the complex commutator form is
\begin{equation}\label{eq:finite-form}
 q_{s,t}(\xi)
      =u\|[\xi,E_{12}]\|_{\HS}^2
           +v\|[\xi,E_{23}]\|_{\HS}^2.
\end{equation}
For the geometric mass on $A\otimes\C^3$,
the eigenvalues below are divided by
$2\Vol(\omega_0)$.

\begin{proposition}\label{prop:finite-spectrum}
The eight eigenvalues of the complex Hermitian
form \eqref{eq:finite-form}, listed without
sorting, are
\begin{equation}\label{eq:eigenvalues}
 \begin{gathered}
 0,\ v,\ u,\ 2u,\ 2v,\ u+v,\\
 u+v-\sqrt{u^2-uv+v^2},\qquad
 u+v+\sqrt{u^2-uv+v^2}.
 \end{gathered}
\end{equation}
Its complex kernel has dimension one when
$uv>0$, four when exactly one weight vanishes,
and eight at the origin.
The real Hermitian restriction is discussed below.
For $s=\varepsilon\delta$, $t=\varepsilon$,
with positive $\varepsilon,\delta\to0$, three
positive eigenvalues have order
$\varepsilon^2\delta^2$ and four have order
$\varepsilon^2$.
\end{proposition}

\begin{proof}
The orthonormal off-diagonal units
$E_{13},E_{12},E_{23},E_{21},E_{32},E_{31}$
are eigenvectors with eigenvalues
$0,v,u,2u,2v,u+v$, respectively. The identity
$[E_{ij},E_{kl}]=\delta_{jk}E_{il}-\delta_{li}E_{kj}$
shows that their images in the direct sum of the
two target matrix spaces are orthogonal to each
other and to the images of diagonal matrices.
On $\xi=\diag(d_1,d_2,d_3)$ with
$d_1+d_2+d_3=0$, the form is
$u|d_1-d_2|^2+v|d_2-d_3|^2$, with mass
$\sum|d_i|^2$. Its eigenvalues on the trace-free
plane are
the two roots other than the constant-vector
root of the weighted path matrix
\[
 \begin{pmatrix}
 u&-u&0\\
 -u&u+v&-v\\
 0&-v&v
 \end{pmatrix}.
\]
Their sum is $2(u+v)$ and product $3uv$,
which gives \eqref{eq:eigenvalues}, including
the limiting zero weights. For $u/v=\delta^2\to0$,
the smaller root is
\[
 \frac{3uv}{u+v+\sqrt{u^2-uv+v^2}}
       =\tfrac32u(1+O(u/v)),
\]
and the larger is $2v+O(u)$. This proves both
scale counts and the kernel assertions.
\end{proof}

These are complex commutator eigenvalues, not
automatically HYM relaxation rates at a
nonbalanced point. Their kernel dimensions also
differ from those on the real Hermitian algebra.
For $uv>0$ the complex kernel is $\C E_{13}$
and contains no nonzero Hermitian element.
On a nonzero one-arrow axis the trace-free
Hermitian kernel has real dimension one;
at the origin it has dimension eight.
Thus the complex dimensions are $1,4,8$,
whereas the real Hermitian dimensions are $0,1,8$.

The same fixed smooth background therefore
contains actual noncommuting integrable points
with changing stabilizers and arbitrarily
separated positive finite scales. The general
theorem retains every finite direction throughout
this family. For any associated HYM trajectory
satisfying its confinement hypothesis, the
selected exact graph flow, $C^1$ initial phase
and exponential relative error have the same
uniform constants. Determining confinement
and the algebraic asymptotics of that corrected
finite motion are further dynamical questions.
